\documentclass{article}
\usepackage[T1]{fontenc}
\usepackage{lmodern}
\usepackage{graphicx}
\usepackage{amsmath,amssymb}
\usepackage[a4paper,margin=2cm]{geometry}
\usepackage[absolute,overlay]{textpos}
\setlength{\TPHorizModule}{1mm}
\setlength{\TPVertModule}{1mm}
\usepackage[indonesian]{babel}
\usepackage{hyperref}
\setlength{\emergencystretch}{2em}
% unit: Halaman 61

\begin{document}

% === Halaman 61 (llm) ===

\begin{itemize}
    \item Cipher berulang dinyatakan sebagai
    \[ C_i = f(C_{i-1}, K_i) \]
    yang dalam hal ini,
    \begin{itemize}
        \item $i = 1, 2, ..., r$ ($r$ adalah jumlah putaran).
        \item $K_i$ = kunci putaran (round key) pada putaran ke-$i$
        \item $f$ = fungsi transformasi (round function), di dalamnya terdapat operasi substitusi, permutasi, dan/atau ekspansi, kompresi).
    \end{itemize}
    \item Plainteks dinyatakan dengan $C_0$ dan cipherteks dinyatakan dengan $C_r$.
\end{itemize}

\end{document}
